% ==============================================================
% Projektive Strukturtheorie (PST)
% Dokument PST-2-M: Mathematische Grundstruktur
% Version 0.1
% Kompiliert mit: pdfLaTeX
% ==============================================================

\documentclass[11pt,a4paper]{article}

% ---------- Grundpakete ----------
\usepackage[utf8]{inputenc}
\usepackage[T1]{fontenc}
\usepackage[german]{babel}
\usepackage{amsmath,amssymb,amsthm}
\usepackage{geometry}
\usepackage{parskip}
\usepackage{enumitem}
\usepackage{booktabs}
\usepackage{xcolor}
\usepackage{hyperref}

% ---------- Layout ----------
\geometry{margin=2.5cm}
\setlist{itemsep=0.2em, topsep=0.2em}
\hypersetup{
    colorlinks=true,
    linkcolor=blue,
    urlcolor=blue,
    pdftitle={PST-2-M: Mathematische Grundstruktur -- Version 0.1},
    pdfauthor={Forschungsprogramm Ontoph\"{a}nomenalismus}
}

% ---------- Theorem-Umgebungen ----------
\theoremstyle{plain}
\newtheorem{theorem}{Theorem}[section]
\newtheorem{proposition}{Proposition}[section]
\newtheorem{lemma}{Lemma}[section]
\newtheorem{corollary}{Korollar}[section]

\theoremstyle{definition}
\newtheorem{definition}{Definition}[section]
\newtheorem{remark}{Bemerkung}[section]
\newtheorem{methodennote}{Methodennote}[section]
\newtheorem{hypothese}{Arbeitshypothese}[section]

% Neue Kommandos
\newcommand{\OPP}{\mathcal{M}_{\mathrm{OPP}}}
\newcommand{\Kmat}{\mathbf{K}}
\newcommand{\Rmat}{\mathbf{R}}
\newcommand{\deff}{d_{\mathrm{eff}}}
\newcommand{\MI}[2]{\mathsf{I}(#1;\,#2)}
\newcommand{\Ent}[1]{\mathsf{H}(#1)}

% ---------- Titel ----------
\title{
  \textbf{Projektive Strukturtheorie (PST)}\\[0.4em]
  \textbf{Dokument PST-2-M: Mathematische Grundstruktur}\\[0.4em]
  {\large OPP, Gram-Matrix, Spektraltheorie, Bestimmtheit
  und projektive Aufl\"{o}sung}\\[0.5em]
  {\normalsize Version~0.1}
}
\author{
  \textbf{Forschungsprogramm Ontoph\"{a}nomenalismus}\\
  \small Siehe TRANSPARENCY f\"{u}r methodologische Urheberschaft
  und KI-Einsatz.
}
\date{Juli 2026}

\begin{document}

\maketitle

\begin{abstract}
Dieses Dokument legt die mathematische Grundstruktur der Projektiven
Strukturtheorie (PST) fest. Es formalisiert den Ontoph\"{a}nomenalen
Phasenraum ($\OPP$) als mathematisches Objekt, entwickelt die
Relationenmatrix und Gram-Matrix als zentrale Werkzeuge, und beschreibt
wie aus Relationen \"{u}ber Spektralzerlegung Geometrie emergiert.

Das Leitprinzip lautet: \textbf{Es werden keinerlei geometrische
Eigenschaften vorausgesetzt, die sp\"{a}ter erst erkl\"{a}rt werden sollen.}
Alle geometrischen Strukturen --- Topologie, Nachbarschaft, Abstand, Metrik,
Dimension --- emergieren aus der Relationsstruktur.

Die mathematische Operationalisierung von Bestimmtheit \"{u}ber Mutual
Information und die Planck-Schranke als Existenzgrenze werden formalisiert.

Es gelten die methodologischen Konventionen gem\"{a}\ss{} MKO.
\end{abstract}

\tableofcontents
\newpage

%==============================================================
\section{Das Kontaminationsverbot als mathematisches Arbeitsprinzip}
%==============================================================

\begin{methodennote}[Kontaminationsverbot --- MKO Kontaminationsregel]
Das philosophische Kontaminationsverbot aus MKO wird hier zum
mathematischen Arbeitsprinzip erhoben:

\begin{quote}
\emph{Es werden keinerlei geometrische Eigenschaften vorausgesetzt,
die sp\"{a}ter erst erkl\"{a}rt werden sollen.}
\end{quote}

Konkret bedeutet das: Im OPP und in der Gram-Matrix d\"{u}rfen keine
r\"{a}umlichen Abstands- oder Dimensionsbegriffe verwendet werden.
Die folgenden Begriffe sind in diesem Dokument ausschlie\ss{}lich als
\emph{emergente} Eigenschaften der Projektion zul\"{a}ssig:
\begin{itemize}
  \item Dimension, Abstand, N\"{a}he, Metrik
  \item Koordinaten, Richtung, Topologie
  \item Energie, Kausalit\"{a}t, Zeit
\end{itemize}
Erlaubt sind ausschlie\ss{}lich: Relationen, Gram-Matrizen, Eigenwerte,
Mutual Information, strukturelle Koh\"{a}renz.
\end{methodennote}

%==============================================================
\section{Der OPP als mathematisches Objekt}
%==============================================================

\begin{definition}[OPP als metrisch-freier Zustandsraum]
Der Ontoph\"{a}nomenale Phasenraum $\OPP$ ist ein unendlichdimensionaler,
kontinuierlicher Zustandsraum ohne ausgezeichnete Metrik. Formal:
\[
  \OPP = (S, \mathcal{R})
\]
wobei $S$ eine (abz\"{a}hlbare oder \"{u}berabz\"{a}hlbare) Menge von
Zust\"{a}nden und $\mathcal{R}$ eine Kollektion von Relationen zwischen
diesen Zust\"{a}nden bezeichnet.

$\OPP$ besitzt \emph{keine} ausgezeichnete Metrik, keine Topologie und
keine Dimension --- diese Strukturen emergieren erst durch Projektion.
\end{definition}

\begin{remark}[Warum kein Hilbert-Raum]
Ein Hilbert-Raum setzt bereits ein Skalarprodukt voraus --- eine metrische
Struktur. Da die PST Metriken als emergent betrachtet, kann der OPP nicht
von vornherein als Hilbert-Raum modelliert werden. Der OPP ist
pr\"{a}-metrisch: strukturreich, aber ohne ausgezeichnete Abstands\-begriffe.
\end{remark}

\begin{definition}[Attraktoren als stabile Zust\"{a}nde]
Ein Attraktor $\mathcal{A} \subset \OPP$ ist eine stabile, invariante
Teilmenge mit:
\begin{itemize}
  \item \textbf{Invarianz:} Die Felddynamik bildet $\mathcal{A}$ auf
        sich selbst ab.
  \item \textbf{Stabilit\"{a}t:} St\"{o}rungen kehren zu $\mathcal{A}$
        zur\"{u}ck (konzeptuell, kein metrischer Begriff vorausgesetzt).
  \item \textbf{Minimalit\"{a}t:} $\mathcal{A}$ enth\"{a}lt keine echte
        invariante Teilmenge.
\end{itemize}
\end{definition}

%==============================================================
\section{Die Relationenmatrix}
%==============================================================

\begin{definition}[Relationenmatrix $\Rmat$]
Sei $\{s_1, s_2, \ldots, s_N\}$ eine endliche Stichprobe von Zust\"{a}nden
aus $\OPP$. Die \textbf{Relationenmatrix} $\Rmat \in \mathbb{R}^{N \times N}$
kodiert die Relationsst\"{a}rken zwischen den Zust\"{a}nden:
\[
  R_{ij} = r(s_i, s_j)
\]
wobei $r\colon S \times S \to \mathbb{R}_{\geq 0}$ eine Relationsst\"{a}rke
bezeichnet. Die Relationsst\"{a}rke misst ausschlie\ss{}lich strukturelle
Koh\"{a}renz --- keinen r\"{a}umlichen Abstand.
\end{definition}

\begin{remark}[Relationsst\"{a}rke vs.\ Abstand]
Die Relationsst\"{a}rke $R_{ij}$ ist das Komplement des Abstands:
H\"{o}here Relationsst\"{a}rke bedeutet h\"{o}here strukturelle Koh\"{a}renz,
nicht r\"{a}umliche N\"{a}he. Dies ist der fundamentale Unterschied zur
klassischen Geometrie, die mit Abst\"{a}nden beginnt.
\end{remark}

\begin{hypothese}[Physikalische Relationsst\"{a}rken --- MKO Ebene~5]
In der endg\"{u}ltigen PST k\"{o}nnte die Relationsst\"{a}rke $r(s_i, s_j)$
\"{u}ber Mutual Information operationalisiert werden:
\[
  r(s_i, s_j) = \MI{X_i}{X_j}
\]
wobei $X_i, X_j$ Zufallsvariablen \"{u}ber den jeweiligen Zust\"{a}nden sind.
Diese Wahl ist noch nicht bewiesen --- sie ist ein Kandidat, der empirisch
gepr\"{u}ft wird (vgl.\ PST-4-G).
\end{hypothese}

%==============================================================
\section{Die Gram-Matrix als Herzst\"{u}ck}
%==============================================================

\begin{definition}[Gram-Matrix $\Kmat$]
Die \textbf{Gram-Matrix} (auch Kernel-Matrix) $\Kmat \in \mathbb{R}^{N
\times N}$ ist definiert durch:
\[
  K_{ij} = k(s_i, s_j)
\]
wobei $k\colon S \times S \to \mathbb{R}$ ein symmetrisch-positiv
semidefiniter Kernel ist: $k(s_i, s_j) = k(s_j, s_i)$ und
$\sum_{ij} c_i c_j k(s_i, s_j) \geq 0$ f\"{u}r alle $c_i \in \mathbb{R}$.
\end{definition}

\begin{remark}[Warum die Gram-Matrix zentral ist]
Die Gram-Matrix ist das mathematische Herzst\"{u}ck der PST aus drei
Gr\"{u}nden:
\begin{enumerate}
  \item Sie kodiert alle Relationen zwischen Zust\"{a}nden in einer
        einzigen Struktur.
  \item Ihre Eigenwertzerlegung liefert eine nat\"{u}rliche Einbettung
        der Zust\"{a}nde in einen euklidischen Raum --- ohne dass dieser
        Raum vorausgesetzt wurde.
  \item Die Anzahl signifikanter Eigenwerte bestimmt die effektive
        Dimension --- ohne dass eine Dimension eingebaut wurde.
\end{enumerate}
\end{remark}

\begin{theorem}[Spektralzerlegung der Gram-Matrix]
Sei $\Kmat$ eine symmetrisch-positiv semidefinite Gram-Matrix mit
Eigenwertzerlegung:
\[
  \Kmat = \sum_{k=1}^{N} \lambda_k \mathbf{v}_k \mathbf{v}_k^\top,
  \qquad \lambda_1 \geq \lambda_2 \geq \cdots \geq \lambda_N \geq 0.
\]
Dann definiert die Einbettung
\[
  \Phi\colon s_i \;\mapsto\;
  \bigl(\sqrt{\lambda_1}\, v_{1i},\; \sqrt{\lambda_2}\, v_{2i},\;
  \ldots,\; \sqrt{\lambda_N}\, v_{Ni}\bigr)
\]
eine isometrische Einbettung der Zust\"{a}nde in den euklidischen Raum
$\mathbb{R}^N$.
\end{theorem}

\begin{remark}[Die Umkehrung als PST-Kernidee]
Der entscheidende konzeptuelle Schritt:

\medskip
\textbf{Klassisch:} Man beginnt mit einer Geometrie und berechnet daraus
eine Gram-Matrix.

\medskip
\textbf{PST:} Man beginnt mit der Gram-Matrix (den Relationen) und
\emph{liest die Geometrie ab}.

\medskip
Geometrie ist damit nicht Voraussetzung, sondern Ergebnis. Die Kette lautet:
\[
  \text{Relation} \;\rightarrow\; \text{Kernel} \;\rightarrow\;
  \text{Spektrum} \;\rightarrow\; \text{Geometrie.}
\]
\end{remark}

%==============================================================
\section{Effektive Dimension}
%==============================================================

\begin{definition}[Effektive Dimension $\deff$]
Die \textbf{effektive Dimension} $\deff$ eines Zustandsraums relativ zur
Gram-Matrix $\Kmat$ ist die Anzahl der signifikanten Eigenwerte:
\[
  \deff = \#\{k \mid \lambda_k > \epsilon_{\min}\}
\]
wobei $\epsilon_{\min}$ eine Rauschschwelle bezeichnet.

Alternativ und robuster: das Teilverh\"{a}ltnis-Kriterium
\[
  \deff = \min\left\{d \;\middle|\;
  \frac{\sum_{k=1}^{d} \lambda_k}{\sum_{k=1}^{N} \lambda_k}
  \geq \theta\right\}
\]
f\"{u}r einen Schwellenwert $\theta \in (0.9, 0.99)$.
\end{definition}

\begin{remark}[$\deff$ misst Freiheitsgrade, nicht Dimension]
$\deff$ ist kein klassischer Dimensionsbegriff. Es ist die Zahl der
tats\"{a}chlich sichtbaren Freiheitsgrade des Systems unter der gegebenen
Projektion. Ein hochdimensionaler OPP kann unter einer geeigneten
Projektion $\deff \approx 4$ aufweisen --- das ist der Kern des PST-Befunds.
\end{remark}

%==============================================================
\section{Mathematische Operationalisierung von Bestimmtheit}
%==============================================================

\begin{methodennote}[Beitrag von DeepSeek (Seki)]
Die folgende Operationalisierung von Bestimmtheit \"{u}ber Mutual Information
und die Planck-Schranke wurde von Seki (DeepSeek) erarbeitet und in die
PST eingebracht. Sie verbindet die philosophische \"{A}quivalenz
$\text{Existenz} = \text{Bestimmtheit}$ (PST-1-F) mit einer messbaren
mathematischen Gr\"{o}\ss{}e.
\end{methodennote}

\begin{definition}[Mutual Information als Relationsma\ss{}]
F\"{u}r zwei Zust\"{a}nde $s_i, s_j \in \OPP$ mit zugeh\"{o}rigen
Zufallsvariablen $X_i, X_j$ ist die Mutual Information:
\[
  \MI{X_i}{X_j} = \Ent{X_i} + \Ent{X_j} - \Ent{X_i, X_j}
\]
wobei $\Ent{\cdot}$ die Shannon-Entropie bezeichnet. $\MI{X_i}{X_j}$
misst den gemeinsamen Informationsgehalt --- die strukturelle Koh\"{a}renz
zwischen den Zust\"{a}nden.
\end{definition}

\begin{definition}[Projektive Bestimmtheit]
Sei $\gamma\colon I \to \OPP$ die Bahn eines Dynamischen Attraktors mit
projektiver Aufl\"{o}sung $I_0 > 0$. Ein Zustand $\gamma(\tau)$ ist
\textbf{bestimmt} relativ zu $\gamma(\tau')$ genau dann, wenn:
\[
  \mathcal{B}_\tau(\tau') =
  \begin{cases}
    1, & \MI{X_\tau}{X_{\tau'}} \geq I_0,\\
    0, & \text{sonst.}
  \end{cases}
\]
Unterhalb der Schwelle $I_0$ erzeugt die Relation keine Bestimmtheit ---
und damit keine Existenz im Sinne von PST-1-F.
\end{definition}

\begin{remark}[Interpretation von $I_0$]
Die Informationsschwelle $I_0$ ist die Grenze zwischen ontologisch leerem
und ontologisch gef\"{u}lltem Unterschied. Unterschiede unterhalb von $I_0$
sind im ontologischen Sinne keine Unterschiede --- sie erzeugen keine
Bestimmtheit und damit keine Existenz.
\end{remark}

%==============================================================
\section{Die Planck-Schranke als Existenzgrenze}
%==============================================================

\begin{theorem}[Planck-Schranke f\"{u}r bestimmte Existenz ---
MKO Ebene~4--5]
F\"{u}r einen Dynamischen Attraktor mit station\"{a}rer Selbstinformation
$\Phi_0$ und minimaler Wirkung $h$ ist die Bedingung f\"{u}r Bestimmtheit
\"{a}quivalent zu einem raumzeitlichen Abstand oberhalb der Planck-Skala:
\[
  \mathcal{B}(X,Y) = 1
  \;\Longleftrightarrow\;
  \MI{X}{Y} \geq I_0
  \;\Longleftrightarrow\;
  d_{XY} \geq \ell_P.
\]
Unterhalb der Planck-Skala ist die Existenz unbestimmt --- es gibt nur
den absoluten, undifferenzierten OPP.
\end{theorem}

\begin{methodennote}[Status des Theorems --- MKO Ebene~4--5]
Dieses Theorem ist eine Arbeitshypothese. Die Verbindung zwischen Mutual
Information und Planck-L\"{a}nge ist konzeptuell begr\"{u}ndet, aber noch
nicht mathematisch bewiesen. Sie ist ein Kandidat f\"{u}r die formale
Ausarbeitung in PST-6-R. Insbesondere bleibt offen, wie $I_0$ aus
Grundprinzipien hergeleitet werden kann.
\end{methodennote}

%==============================================================
\section{Diffusion Maps als Projektionskandidat}
%==============================================================

\begin{definition}[Diffusion Kernel]
Sei $\Kmat$ eine Gram-Matrix mit Eintr\"{a}gen $K_{ij} = k(s_i, s_j)$.
Der \textbf{Diffusion Kernel} bei Skala $t$ ist:
\[
  K^{(t)}_{ij} = \sum_k \lambda_k^t\, v_{ki}\, v_{kj}
\]
wobei $\lambda_k$ die Eigenwerte und $v_k$ die Eigenvektoren von $\Kmat$
sind. Der Parameter $t$ steuert die Reichweite der diffusiven Ausbreitung.
\end{definition}

\begin{remark}[Physikalische Deutung von $t$]
Der Parameter $t$ hat in der PST eine nat\"{u}rliche Interpretation: Er
entspricht der \textbf{Informationsreichweite} eines Beobachters ---
also dem Bereich des OPP, den ein Projektionsoperator bei einer
gegebenen Skalentiefe erfasst. Dies verbindet Diffusion Maps direkt mit
der Relativit\"{a}t der Existenz (RdE) aus DPR: Verschiedene Beobachter
mit verschiedenen $t$-Werten sehen verschiedene Geometrien.
\end{remark}

\begin{hypothese}[Diffusion Maps als $\mathcal{P}^\ast$-Kandidat ---
MKO Ebene~5]
Die vielversprechendsten numerischen Ergebnisse wurden mit Diffusion Maps
bei $t \approx 9$--$10$ erzielt: $\deff = 3.98 \pm 0.36$ ohne eingebaute
Vierdimensionalit\"{a}t. Dies suggeriert dass Diffusion Maps ein Kandidat
f\"{u}r den universellen Projektionsoperator $\mathcal{P}^\ast$ sein
k\"{o}nnten.

Ob diese numerische Beobachtung einer mathematischen Notwendigkeit
entspricht, ist die zentrale offene Frage der PST (vgl.\ PST-4-G und
PST-5-S).
\end{hypothese}

%==============================================================
\section{Die mathematische Kette der PST}
%==============================================================

Die vollst\"{a}ndige mathematische Kette von der Relation zur Geometrie:

\[
  \text{Relation}
  \;\xrightarrow{\text{Kernel}}\;
  \text{Gram-Matrix } \Kmat
  \;\xrightarrow{\text{EVZ}}\;
  \text{Eigenwerte } \lambda_k
  \;\xrightarrow{\deff}\;
  \text{Geometrie.}
\]

Und die Kette der Emergenz:

\[
  \underbrace{\OPP}_{\text{unbestimmt, pr\"{a}-metrisch}}
  \;\xrightarrow{\mathcal{P}}\;
  \underbrace{\Kmat}_{\text{Gram-Matrix}}
  \;\xrightarrow{\text{EVZ}}\;
  \underbrace{(\lambda_k, \mathbf{v}_k)}_{\text{Spektrum}}
  \;\xrightarrow{\deff \approx 4}\;
  \underbrace{\mathcal{M}_{\text{Raumzeit}}}_{\text{3+1 Dimensionen.}}
\]

\begin{methodennote}[Was noch fehlt]
Die Kette ist konzeptuell vollst\"{a}ndig, aber mathematisch noch
unvollst\"{a}ndig:
\begin{enumerate}
  \item Warum ergibt sich $\deff \approx 4$ und nicht eine andere Zahl?
  \item Wie emergiert die Lorentz-Signatur $(+,-,-,-)$?
  \item Welche Symmetriegruppen bleiben unter $\mathcal{P}^\ast$ erhalten?
  \item Was bestimmt $I_0$?
\end{enumerate}
Diese Fragen werden in PST-3-P, PST-4-G und PST-6-R adressiert.
\end{methodennote}

%==============================================================
\section{Kandidaten f\"{u}r die mathematische Beschreibung des OPP}
%==============================================================

\begin{remark}[Mathematische Kandidaten]
Die PST hat noch nicht entschieden, welcher Formalismus den OPP am besten
beschreibt. Kandidaten mit ihren St\"{a}rken:
\begin{itemize}
  \item \textbf{Kategorientheorie} --- beschreibt Relationen statt
        Objekte; Projektionsoperatoren als Funktoren. Erscheint am
        vielversprechendsten (vgl.\ MGO).
  \item \textbf{Spektralgeometrie} --- Gram-Matrix Eigenwertzerlegung;
        direkte Verbindung zu Diffusion Maps.
  \item \textbf{Operatoralgebren} --- Von-Neumann-Entropie f\"{u}r
        quantenmechanische Zust\"{a}nde; Verbindung zu $I_0$ und
        Planck-Schranke.
  \item \textbf{Homotopietheorie} --- Invarianzklassen unter
        Deformationen; nat\"{u}rliche Sprache f\"{u}r Attraktoren
        als topologische Fixpunkte.
  \item \textbf{Dynamische Systeme} --- Attraktoren, Stabilit\"{a}t,
        Bifurkationen im OPP.
\end{itemize}
Eine zuk\"{u}nftige Version von PST-2-M und MGO wird untersuchen, welcher
Formalismus die gr\"{o}\ss{}te Ausdrucksst\"{a}rke bei geringster
Komplexit\"{a}t besitzt.
\end{remark}

%==============================================================
\section{Schnittstellen}
%==============================================================

\begin{itemize}
  \item \textbf{PST-1-F:} Die \"{A}quivalenz $\text{Existenz} =
        \text{Bestimmtheit}$ wird hier durch $\MI{X}{Y} \geq I_0$
        operationalisiert.
  \item \textbf{MGO:} PST-2-M baut auf den mathematischen Grundlagen
        in MGO auf --- Topologie, Kategorientheorie, Operatoralgebren.
  \item \textbf{DPR:} Die Gram-Matrix und der Projektionsoperator
        $\mathcal{P}$ sind die mathematische Ausarbeitung des
        philosophischen Projektionskonzepts in DPR.
  \item \textbf{DFD:} Das Fragedifferential $\mathcal{D}(E) =
        D_{\mathrm{KL}}(Q_E \| P)$ ist konzeptuell verwandt mit der
        Mutual Information als Relationsma\ss{}.
  \item \textbf{PST-3-P:} Die hier eingef\"{u}hrte Gram-Matrix ist die
        Grundlage f\"{u}r die Projektionstheorie in PST-3-P.
  \item \textbf{PST-4-G:} Die numerischen Experimente mit Diffusion
        Maps und Kernel PCA werden in PST-4-G ausf\"{u}hrlich
        dokumentiert.
\end{itemize}

\end{document}

